% Cover staff member is listed for the FAF-241 TMPS seminar in the FCIM timetable.
\documentclass{utmreport}
\selectlanguage{english}

\institution{Technical University of Moldova}
\faculty{Faculty of Computers, Informatics, and Microelectronics}
\department{Department of Software Engineering}
\studyprogramme{Software Engineering Study Programme}
\documenttype{Laboratory Report No. 1}
\reporttitle{Three Creational Patterns in a Task Reporting Project}
\student{Lucian-Adrian Gavril}{FAF-241}
\supervisor{Daniel Istrati}{asist. univ.}
\cityyear{Chișinău}{2026}
\utmlogo{utm-logo.jpeg}

\begin{document}
\makeutmreporttitlepage
\tableofcontents

\begin{utmabstract}[Abstract]
This report explains the Builder, Factory Method, and Abstract Factory patterns in a task-reporting program. The builder collects a title and minimum priority into one options value. The program filters a task list, then uses a report creator to format the selected tasks as CSV. A second component uses an abstract factory to create a matching Markdown heading and task-row formatter. The report traces the input through these components, maps each pattern to its implementation, and gives build commands and oral-review questions. The same project includes the behavioral continuation in Laboratory 3 and uses one command-line client for both demonstrations.

\utmkeywords{creational patterns, builder, factory method, abstract factory, task report}
\end{utmabstract}

\utmintroduction
The TMPS Laboratory 1 assignment asks for three creational patterns from its list and familiarity with the others. It marks Object Pooling as useful to know rather than required \cite{brief,repo}. This project selects Builder, Factory Method, and Abstract Factory.

The example is a small task-reporting program. Its objectives are:
\begin{utmenumerate}
\item explain how the builder turns a sequence of settings into one options value;
\item trace how a creator chooses a CSV formatter and how a theme factory creates a matching Markdown header and row;
\item follow the three input tasks through the filters and connect the result to the client and shared model in Listings~\ref{lst:client} and~\ref{lst:model}.
\end{utmenumerate}

\chapter{Selected patterns and implementation}
The course's recommended references include the standard design-pattern catalog by Gamma, Helm, Johnson, and Vlissides \cite{patterns}. This chapter maps three choices from the assignment list to their concrete roles in the client and pattern files shown in Listings~\ref{lst:client} and~\ref{lst:creational}.

\section{Builder: collect report options}
\texttt{TaskReportOptionsBuilder} holds a title, an include-completed flag, and a minimum priority. Each setter changes one setting and returns the builder, so the calls can be chained. \texttt{Build} copies the chosen values into one \texttt{TaskReportOptions} record.

In \texttt{RunLab01}, the program sets the title to \texttt{Study tasks} and the minimum priority to 2. The option to include completed tasks remains false, which is the builder's default. The result is a single configuration value passed to the task selector and the output components.

This keeps several options together instead of passing separate parameters through each call. The example does not validate the title or priority in \texttt{Build}; a production version would need to decide what values are valid.

\section{Factory Method: choose one report formatter}
\texttt{TaskReportCreator.Build} asks its subclass for a formatter, then delegates the output.

\texttt{CsvTaskReportCreator} creates the CSV formatter.

\texttt{TextTaskReportCreator} creates the plain-text formatter.

\texttt{RunLab01} constructs the CSV creator. The same build workflow can use a different formatter by choosing another creator subtype. The pattern is useful here because the creation decision belongs to the creator subclass, while formatting behavior stays in the formatter.

\section{Abstract Factory: create a matching board style}
\texttt{IBoardThemeFactory} creates two related products: an \texttt{IBoardHeader} and an \texttt{IBoardTaskRow}. \texttt{MarkdownBoardThemeFactory} returns Markdown versions of both. \texttt{PlainTextBoardThemeFactory} returns plain-text versions.

\texttt{TaskBoardRenderer} receives the factory, stores the two products, and uses them to render a title and rows. The program selects the Markdown factory, keeping the heading and row formats consistent. The abstract factory is useful because the renderer depends on interfaces and does not construct a mismatched header and row itself. Laboratory 3 reuses this renderer in its behavioral client.

\chapter{Source-level trace and results}
This chapter follows the exact task values and calls in \texttt{Program.RunLab01}. It is a source-level trace rather than a performance test. The complete client, shared task model, and creational types appear in Listings~\ref{lst:client}, \ref{lst:model}, and~\ref{lst:creational}.

\section{Filter the input tasks}
The program starts with three shared \texttt{TaskItem} objects. The first filter reads \texttt{IncludeCompleted}, which remains false. It removes task 3 because that task is marked done.

The second filter keeps tasks whose priority is at least 2. The result preserves the original order of the two selected records.

\begin{table}[htbp]
\caption{Input tasks and selection result}\label{tab:selected}
\centering
\begin{tabularx}{0.94\textwidth}{|c|Y|c|c|c|}\hline
\textbf{ID} & \textbf{Task} & \textbf{Priority} & \textbf{Initial state} & \textbf{Selected}\\\hline
1 & Write report & 3 & Open & Yes\\\hline
2 & Review notes & 2 & Open & Yes\\\hline
3 & Archive old draft & 1 & Done & No\\\hline
\end{tabularx}
\end{table}

Task 3 is excluded by the completed-task filter. The minimum-priority filter then retains task 1 at priority 3 and task 2 at priority 2. No sorting call is made, so the output keeps their initial order.

\section{Create the two outputs}
The CSV creator receives the title and selected tasks. Its formatter writes the title, a header row, then one row per selected task. The theme factory independently supplies the Markdown header and task-row implementations used by \texttt{TaskBoardRenderer}. The code therefore demonstrates one Factory Method decision for one formatter and one Abstract Factory decision for a family of two related rendering products.

The source predicts these outputs:
\begin{lstlisting}
Study tasks
Id,Title,Priority,Status
1,"Write report",3,Open
2,"Review notes",2,Open

# Study tasks
- [ ] Write report (P3)
- [ ] Review notes (P2)
\end{lstlisting}

The complete build and run commands appear in the project README and are repeated below for convenient review. The source-level trace does not claim benchmark results.

\clearpage
\section{Build and run}
From the \texttt{lab-01} directory, use the .NET 10 SDK. Build artifacts are directed outside the project folder:
\begin{lstlisting}
$build = Join-Path $env:TEMP 'tmps-build'
dotnet build .\TMPSLab01.csproj -c Release --artifacts-path $build
dotnet (Join-Path $build 'bin\TMPSLab01\release\TMPSLab01.dll') lab-01
\end{lstlisting}

The build command compiles both laboratory modules into one executable. The \texttt{lab-01} argument selects this report's demonstration. The program prints the selected tasks first in CSV form and then in Markdown form. The plain-text report creator and plain-text board factory are implemented but are not selected by \texttt{RunLab01}.

\section{Questions for oral review}
\begin{enumerate}
\item What values does \texttt{TaskReportOptionsBuilder} collect? It stores a title, whether completed tasks are included, and a minimum priority.
\item What decision does Factory Method make in this example? A creator subclass chooses the formatter returned by its creation method.
\item What makes the theme factory abstract? It supplies a matching header and task-row pair through interfaces.
\item How can the program include completed tasks? Call \texttt{IncludeCompletedTasks} while building the options.
\item How would a JSON report fit the existing design? Add a JSON formatter and a creator subtype that returns it.
\end{enumerate}

\utmconclusions
The example applies three creational patterns to separate decisions about configuration and object creation. Builder assembles one options value. Factory Method chooses the report formatter. Abstract Factory creates a matching pair of board renderers. The source trace shows that the selected settings remove the completed low-priority task and preserve the order of the two open tasks.

The example has limits: the builder does not validate settings, the sample uses in-memory data, and this mode exercises only one creator and one theme factory at runtime. The shared executable also contains the behavioral continuation, which uses the same task model and creational components.

\begin{thebibliography}{9}
\bibitem{brief} Technical University of Moldova, ELSE. \textit{FAF.TMPS21.1: Laboratory work \#1}, Creational Design Patterns. Available online: \url{https://else.fcim.utm.md/mod/assign/view.php?id=27825} (accessed 10 October 2026).
\bibitem{repo} V. Drumea. \textit{SDTM-Labs, Lab \#1 README}. Project organization, explanation, and repository submission requirements. Available online: \url{https://github.com/DrVasile/SDTM-Labs/blob/master/Lab%231/README.md} (accessed 10 October 2026).
\bibitem{patterns} E. Gamma, R. Helm, R. Johnson, and J. Vlissides. \textit{Design Patterns: Elements of Reusable Object-Oriented Software}. Addison-Wesley, 1994. The course lists this work in its references section.
\end{thebibliography}

\appendix
\chapter{Source files}\label{app:source}
The following files show the single command-line client, shared task model, and creational pattern types described in this report. The project also compiles \texttt{BehavioralPatterns.cs} from the Laboratory 3 folder into the same assembly.
\section{Client}
\lstinputlisting[caption={Single client with Lab 1 and Lab 3 modes},label={lst:client}]{Program.cs}
\section{Shared model}
\lstinputlisting[caption={Task model shared by both laboratories},label={lst:model}]{TaskItem.cs}
\section{Creational types}
\lstinputlisting[caption={Builder, Factory Method, and Abstract Factory},label={lst:creational}]{CreationalPatterns.cs}
\end{document}
